\lecture

\begin{guiding}
We finish the diagonal computation: the dual class of the diagonal is a sum over
dual bases, and evaluating it turns $e(TM)$ into the Euler characteristic. Then
we introduce the last family of characteristic classes, the Pontryagin classes
of a real bundle.
\end{guiding}

Throughout $M$ is a closed $n$-manifold with fundamental class $\mu$;
coefficients are a field, always admissible over $\Z_2$ and, when $M$ is
oriented, over any field.

\section{The local computation}

\begin{proof}[Proof of Lemma~\ref{lem:u-over-mu}]
\begin{strategy}
The class $u''/\mu$ lies in $H^0(M)$, so it is determined by its value at each
point $x$. Naturality of the slant product turns that value into the evaluation
of the dual class against the local fundamental class at $x$, and the
fibre-restriction property of the Thom class makes it $1$.
\end{strategy}

\step{1} \emph{Reduction to a point.} Fix $x\in M$ and let
$i_x\colon M\to M\times M$, $y\mapsto(x,y)$. Naturality of the slant product in
the first variable gives that the value of $u''/\mu$ at $x$ is
\[
\big(i_x^*u''\big)/\mu=\big\langle i_x^*u'',\ \mu\big\rangle ,
\]
where we use that $i_x^*u''\in H^n(M)$ and that slanting an $n$-class with the
$n$-cycle $\mu$ is the Kronecker pairing.

\step{2} \emph{Passing to the relative class.} Let
$j_x\colon(M,M\smallsetminus x)\to(M\times M,\ M\times M\smallsetminus\Delta)$ be
$y\mapsto(x,y)$; it is well defined, since $(x,y)\in\Delta$ forces $y=x$. Then
$i_x^*u''=i_M^*j_x^*u'$, where $i_M\colon M\to(M,M\smallsetminus x)$, so with
$\mu_x\in H_n(M,M\smallsetminus x)$ the image of $\mu$,
\[
\big\langle i_x^*u'',\mu\big\rangle
=\big\langle j_x^*u',\ \mu_x\big\rangle .
\]

\step{3} \emph{Identifying $j_x^*u'$.} Under the exponential chart of
Theorem~\ref{thm:tnt}, the map $j_x$ corresponds to the inclusion of the fibre
of $N(\Delta)$ over $(x,x)$:
\[
\big(N_x,\ N_x\smallsetminus0\big)
\xrightarrow{\ \mathrm{Exp}\ }
\big(M,\ M\smallsetminus x\big)
\xrightarrow{\ j_x\ }
\big(M\times M,\ M\times M\smallsetminus\Delta\big),
\]
the first map being a homeomorphism onto a neighbourhood, and the composite
being homotopic to $v\mapsto(x,\exp_x v)$. Since $u'$ comes from the Thom class,
its restriction to a fibre pair is the preferred generator; hence $j_x^*u'$ is
the preferred generator of $H^n(T_xM,T_xM\smallsetminus0)$.

\step{4} \emph{Conclusion.} The correspondence between orientations of $T_xM$
and local orientations of $M$ is exactly the one induced by $\exp$, so $\mu_x$
corresponds to the preferred generator in homology and
$\langle j_x^*u',\mu_x\rangle=1$. Therefore $u''/\mu$ takes the value $1$ at
every $x$, that is, $u''/\mu=1$.
\end{proof}

\section{The duality theorem and the Euler characteristic}

\begin{theorem}[Poincar\'e duality, quoted]\label{thm:duality}
For $M$ closed, and with field coefficients (any field if $M$ is oriented,
$\Z_2$ otherwise), the cup product pairing
$H^i(M)\otimes H^{n-i}(M)\to\Lambda$, $(x,y)\mapsto\langle x\smile y,\mu\rangle$,
is non-degenerate. Equivalently, every homogeneous basis $b_1,\dots,b_r$ of
$H^*(M)$ admits a \emph{dual basis} $b_1^\#,\dots,b_r^\#$ with
\[
\big\langle b_i\smile b_j^{\#},\ \mu\big\rangle=\delta_{ij}.
\]
\end{theorem}

\begin{theorem}\label{claim:diag-class-formula}
In terms of such bases the dual class of the diagonal is
\[
u''=\sum_{i=1}^{r}(-1)^{\dim b_i}\ b_i\times b_i^{\#}.
\]
\end{theorem}

\begin{proof}
\begin{strategy}
Expand $u''$ in the K\"unneth basis with unknown coefficients, then exploit the
two facts proved so far: $u''$ slants to $1$ against $\mu$, and multiplying by
$a\times1$ or $1\times a$ gives the same answer. Evaluating on basis elements
identifies the unknowns as the dual basis.
\end{strategy}

\step{1} \emph{Expansion.} By K\"unneth,
$H^n(M\times M)=\bigoplus_i H^{\dim b_i}(M)\otimes H^{n-\dim b_i}(M)$, so we may
write
\[
u''=\sum_{i=1}^{r} b_i\times c_i,
\qquad \dim b_i+\dim c_i=n .
\]

\step{2} \emph{A formula for every class.} Let $a\in H^*(M)$. Using
Lemma~\ref{lem:slant-identity}, then Lemma~\ref{lem:ax1}, then
Lemma~\ref{lem:u-over-mu},
\[
a=a\smile\big(u''/\mu\big)
=\big((a\times1)\smile u''\big)\big/\mu
=\big((1\times a)\smile u''\big)\big/\mu .
\]

\step{3} \emph{Expanding the right-hand side.} Substituting the expansion of
$u''$ and moving $a$ past $b_i$, which costs the Koszul sign
$(-1)^{\dim a\,\dim b_i}$,
\[
\big((1\times a)\smile u''\big)\big/\mu
=\sum_{i}(-1)^{\dim a\,\dim b_i}\,\big(b_i\times(a\smile c_i)\big)\big/\mu
\]
\[
=\sum_{i}(-1)^{\dim a\,\dim b_i}\,
\big\langle a\smile c_i,\ \mu\big\rangle\ b_i .
\]

\step{4} \emph{Identification.} Take $a=b_j$. Comparing with Step 2, which says
the sum equals $b_j$, and using that the $b_i$ are a basis,
\[
(-1)^{\dim b_j\dim b_i}\big\langle b_j\smile c_i,\mu\big\rangle=\delta_{ij}.
\]
For $i=j$ the sign is $(-1)^{(\dim b_i)^2}=(-1)^{\dim b_i}$, so
$\langle b_i\smile c_i,\mu\rangle=(-1)^{\dim b_i}$, while for $i\neq j$ the
pairing vanishes. Comparing with Theorem~\ref{thm:duality} gives
$c_i=(-1)^{\dim b_i}b_i^{\#}$.
\end{proof}

\begin{theorem}\label{thm:euler-char}
For $M$ closed and oriented,
\[
\big\langle e(TM),\ \mu\big\rangle=\chi(M)=\sum_k(-1)^k\dim H^k(M);
\]
without orientation, $\langle w_n(TM),\mu\rangle\equiv\chi(M)\bmod 2$.
\end{theorem}

\begin{proof}
By Corollary~\ref{cor:diagonal-euler} and
Theorem~\ref{claim:diag-class-formula},
\[
e(TM)=\Delta^*u''=\sum_i(-1)^{\dim b_i}\,\Delta^*\big(b_i\times b_i^\#\big)
=\sum_i(-1)^{\dim b_i}\,b_i\smile b_i^{\#} .
\]
Evaluating on $\mu$ and using $\langle b_i\smile b_i^\#,\mu\rangle=1$,
\[
\big\langle e(TM),\mu\big\rangle=\sum_i(-1)^{\dim b_i}
=\sum_k(-1)^k\dim H^k(M)=\chi(M),
\]
since the $b_i$ of a given dimension $k$ form a basis of $H^k(M)$.
\end{proof}

\begin{example}
$\chi(S^2)=2$, so $e(TS^2)\neq0$: a third proof that $TS^2$ is non-trivial, and
the statement underlying the hairy ball theorem. For odd-dimensional $M$ we have
$\chi(M)=0$, consistent with $2e(\xi)=0$ in odd rank.
\end{example}

\section{Chern classes through the Gysin sequence}

We describe the inductive construction of Chern classes used in
Milnor--Stasheff, an alternative to the Grassmannian route of Lecture 6.

\begin{theorem}[Gysin sequence]\label{thm:gysin}
For an oriented $\R^n$-bundle $\xi$ with total space $E\to B$ there is a long
exact sequence with integer coefficients
\[
\cdots\to H^{i-n}(B)\ \xrightarrow{\ \smile e\ }\ H^{i}(B)
\ \xrightarrow{\ \pi_0^*\ }\ H^{i}(E_0)\ \to\ H^{i-n+1}(B)\to\cdots
\]
\end{theorem}

\begin{proof}
Start from the long exact sequence of the pair $(E,E_0)$ and substitute
$H^{j}(E,E_0)\iso H^{j-n}(B)$, by the Thom isomorphism, and
$H^j(E)\iso H^j(B)$, by homotopy invariance. The map
$H^{j-n}(B)\to H^{j}(B)$ so obtained is computed by
\[
i^*\phi(y)=i^*\big(\pi^*y\smile u\big)=\pi^*y\smile i^*u=\pi^*\big(y\smile e\big),
\]
so under the identification it becomes $\smile e$.
\end{proof}

\begin{construction}[Chern classes, inductively]
Let $E\to B$ be a complex bundle of rank $n$, with its canonical orientation.
Set $c_n(E):=e(E)$. Over the sphere bundle $\pi_0\colon S(E)=E_0\to B$ there is
a tautological complex line sub-bundle $L\subset\pi_0^*E$, spanned at a point
$(x,v)$ by $v$; let $W=L^{\perp}$, of rank $n-1$. In degrees $i<2n-1$ the Gysin
sequence shows $\pi_0^*\colon H^i(B)\to H^i(E_0)$ is an isomorphism, so we may
define
\[
c_k(E):=(\pi_0^*)^{-1}c_k(W),\qquad k<n,
\]
inductively on the rank. The two definitions of $c_k$ agree, as one checks on
sums of line bundles using the splitting principle.
\end{construction}

\begin{proposition}
The total Chern class $c(E)=1+c_1+\dots+c_n$ is invertible in $H^\Pi(B;\Z)$.
\end{proposition}

\begin{proof}
The proof of Proposition~\ref{prop:invertible} works verbatim over $\Z$: the
inverse is constructed degree by degree, each coefficient being determined by
the previous ones.
\end{proof}

\section{Conjugate bundles and Pontryagin classes}

\begin{definition}
For a complex bundle $\xi$, the \emph{conjugate bundle} $\bar\xi$ is the same
real bundle equipped with the opposite complex structure $-J$.
\end{definition}

\begin{proposition}\label{prop:conjugate}
$c_k(\bar\xi)=(-1)^k c_k(\xi)$.
\end{proposition}

\begin{proof}
\begin{strategy}
Check it for line bundles, where conjugation reverses the canonical orientation
and hence the sign of the Euler class, then propagate by the splitting principle
and injectivity of $h_n^*$.
\end{strategy}

\step{1} \emph{Line bundles.} For a complex line bundle $L$, a complex basis $v$
of a fibre gives the oriented real basis $(v,Jv)$; passing to $-J$ replaces it
by $(v,-Jv)$, which has the opposite orientation. Hence
$e(\bar L)=-e(L)$ and, since $c_1=e$ for line bundles by
Theorem~\ref{thm:e-gamma-n}, $c_1(\bar L)=-c_1(L)$.

\step{2} \emph{Universal case.} Over $(\CP^\infty)^n$, conjugating
$\gamma_1\oplus\dots\oplus\gamma_1$ conjugates each summand, so by Step 1 and
the Whitney formula
\[
h_n^*c_k(\bar\gamma_n)=\sigma_k(-a_1,\dots,-a_n)=(-1)^k\sigma_k(a_1,\dots,a_n)
=(-1)^k h_n^*c_k(\gamma_n).
\]
Injectivity of $h_n^*$ gives $c_k(\bar\gamma_n)=(-1)^kc_k(\gamma_n)$.

\step{3} \emph{General case.} If $f$ classifies $\xi$ then it also classifies
$\bar\xi$ as a conjugate, in the sense that $f^*\bar\gamma_n\iso\bar\xi$; apply
naturality.
\end{proof}

\begin{definition}
For a real bundle $\xi$, the \emph{complexification} is
$\xi\otimes\C$, whose underlying real bundle is $\xi\oplus\xi$ with complex
structure $J(x,y)=(-y,x)$.
\end{definition}

\begin{lemma}\label{lem:complexification-selfconj}
$\xi\otimes\C\iso\overline{\xi\otimes\C}$. Consequently the odd Chern classes of
a complexification are $2$-torsion:
$2c_{2i+1}(\xi\otimes\C)=0$.
\end{lemma}

\begin{proof}
The map $(x,y)\mapsto(x,-y)$ is a real bundle isomorphism
$\xi\oplus\xi\to\xi\oplus\xi$, and it is conjugate linear for $J$:
\[
J(x,-y)=(y,x),
\qquad\text{while}\qquad
-\overline{J}(x,y)\ \text{corresponds to}\ (y,x),
\]
so it is an isomorphism of complex bundles from $\xi\otimes\C$ onto its
conjugate. By Proposition~\ref{prop:conjugate},
$c_{2i+1}(\xi\otimes\C)=c_{2i+1}(\overline{\xi\otimes\C})
=-c_{2i+1}(\xi\otimes\C)$.
\end{proof}

\begin{definition}[Pontryagin classes]
For a real bundle $E\to B$ set
\[
p_i(E):=(-1)^i\,c_{2i}\big(E\otimes\C\big)\ \in H^{4i}(B;\Z),
\qquad p_0=1,
\]
and $p(E)=1+p_1(E)+p_2(E)+\cdots$. Naturality
$f^*p_i(E)=p_i(f^*E)$ is inherited from the Chern classes.
\end{definition}

\begin{theorem}\label{thm:pontryagin-whitney}
$2\big(p(E_1\oplus E_2)-p(E_1)p(E_2)\big)=0$; that is, the total Pontryagin
class is multiplicative modulo $2$-torsion.
\end{theorem}

\begin{proof}
Complexification is additive:
$(E_1\oplus E_2)\otimes\C\iso(E_1\otimes\C)\oplus(E_2\otimes\C)$. The Whitney
formula for Chern classes gives
$c\big((E_1\oplus E_2)\otimes\C\big)=c(E_1\otimes\C)\,c(E_2\otimes\C)$. Passing
from $c$ to $p$ discards the odd Chern classes, and every discarded term is
$2$-torsion by Lemma~\ref{lem:complexification-selfconj}; multiplying the
difference by $2$ therefore kills it.
\end{proof}

\begin{corollary}\label{cor:pontryagin-of-complex}
If $E$ is a complex bundle then $E\otimes\C\iso E\oplus\bar E$, hence
\[
1-p_1+p_2-\cdots+(-1)^np_n
=\big(1-c_1+c_2-\cdots\big)\big(1+c_1+c_2+\cdots\big),
\]
so for instance $p_1=c_1^2-2c_2$ and $p_2=c_2^2-2c_1c_3+2c_4$.
\end{corollary}

\begin{proof}
For a complex $E$ with structure $J$, the map
\[
E\otimes\C=E\oplus E\longrightarrow E\oplus\bar E,
\qquad
(x,y)\longmapsto\Big(\tfrac12(x+Jy),\ \tfrac12(x-Jy)\Big),
\]
is a complex isomorphism, as one checks by applying $J$ to both sides. Hence
$c(E\otimes\C)=c(E)\,c(\bar E)$, and
Proposition~\ref{prop:conjugate} turns $c(\bar E)$ into the alternating series.
Comparing coefficients of even degree gives the stated identity.
\end{proof}

\begin{example}
$T(S^n)\oplus\nu(S^n)$ is trivial with $\nu$ trivial, so by
Theorem~\ref{thm:pontryagin-whitney} and the torsion-freeness of $H^*(S^n;\Z)$,
\[
p(TS^n)=1 .
\]
Pontryagin classes do not distinguish spheres from parallelizable manifolds.
\end{example}

In the next lecture we compute $p(T\CP^n)$, which is far from trivial, extract
numerical invariants and put them to work on cobordism questions.
